\documentclass{beamer}

\usepackage[utf8]{inputenc}
\usepackage[T5]{fontenc}
\usepackage[vietnamese]{babel}

\usepackage{amsmath,amssymb}
\usepackage{tikz}
\usepackage{listings}
\usepackage{array}
\usepackage{colortbl}
\usepackage{booktabs}
\usepackage{xcolor}

\usetheme{Madrid}
\usecolortheme{dolphin}
% --- màu ---
\definecolor{myblue}{RGB}{30,144,255}
\definecolor{mygreen}{RGB}{60,179,113}
\definecolor{myred}{RGB}{220,20,60}
\definecolor{mypurple}{RGB}{138,43,226}
\definecolor{myorange}{RGB}{255,140,0}
\definecolor{mycyan}{RGB}{0,206,209}
\definecolor{tablehead}{RGB}{255,140,0}
\definecolor{tableheadtxt}{RGB}{255,255,255}
\definecolor{codebg}{RGB}{248,250,252}
\definecolor{codeid}{RGB}{255,140,0}

\setbeamercolor{structure}{fg=myblue}
\setbeamercolor{frametitle}{fg=white,bg=myblue}
\setbeamerfont{frametitle}{series=\bfseries}

% macros to highlight text
\newcommand{\hlblue}[1]{\textcolor{myblue}{#1}}
\newcommand{\hlgreen}[1]{\textcolor{mygreen}{#1}}
\newcommand{\hlred}[1]{\textcolor{myred}{#1}}
\newcommand{\hlpurple}[1]{\textcolor{mypurple}{#1}}
\newcommand{\hlorange}[1]{\textcolor{myorange}{#1}}
\newcommand{\hlcyan}[1]{\textcolor{mycyan}{#1}}

% listings (colorful)
\lstset{
 basicstyle=\ttfamily\small,
 backgroundcolor=\color{codebg},
 keywordstyle=\color{myblue}\bfseries,
 commentstyle=\color{mygreen}\itshape,
 stringstyle=\color{mypurple},
 identifierstyle=\color{codeid},
 numberstyle=\tiny\color{gray},
 breaklines=true,
 showstringspaces=false,
 numbers=none,
 frame=single,
 rulecolor=\color{gray!40},
 emph={int,bool,void,struct,return,vector,queue,priority_queue,auto},
 emphstyle={\color{myorange}\bfseries},
 emph={[2]if,for,while,else},
 emphstyle={[2]\color{myblue}\bfseries},
 emph={[3]find,unite,dfs,bfs,dijkstra},
 emphstyle={[3]\color{mypurple}\bfseries}
}

\title{\hlblue{Lý thuyết đồ thị}}
\author{Dương Khánh Kỳ}
\date{\hlcyan{06/03/2026}}

\begin{document}

\frame{\titlepage}

% 1 Definition
\begin{frame}{\hlorange{Định nghĩa đồ thị}}
\begin{block}{\hlgreen{Khái niệm}}
Đồ thị là mô hình biểu diễn các đối tượng (\hlblue{đỉnh}) và mối quan hệ (\hlred{cạnh}) giữa chúng.
\end{block}

\[
\hlpurple{G}=(\hlblue{V},\hlred{E})
\]

\medskip
\begin{itemize}
\item \hlblue{$V$} — tập các \hlblue{đỉnh} (vertices).
\item \hlred{$E$} — tập các \hlred{cạnh} (edges). Cạnh có thể có trọng số \hlorange{$w$}.
\end{itemize}
\end{frame}

% 2 types: undirected
\begin{frame}{\hlorange{Đồ thị vô hướng (Undirected graph)}}
\begin{block}{\hlgreen{Định nghĩa}}
Trong đồ thị \hlred{vô hướng}, mỗi cạnh không mang hướng:
\[
(\hlblue{u},\hlblue{v}) = (\hlblue{v},\hlblue{u}).
\]
\end{block}

\medskip
\begin{itemize}
\item Thường dùng để mô tả mối quan hệ hai chiều (bạn bè, kết nối mạng).
\item \hlcyan{Bậc (degree)} = số cạnh nối vào đỉnh.
\end{itemize}

\bigskip
\begin{center}
\begin{tikzpicture}[scale=0.9]
\tikzset{v/.style={circle,draw=mygreen,fill=mygreen!20,minimum size=8mm}}
\node[v] (1) at (0,0) {\hlblue{1}};
\node[v] (2) at (2,1) {\hlblue{2}};
\node[v] (3) at (2,-1) {\hlblue{3}};
\draw (1)--(2); \draw (1)--(3); \draw (2)--(3);
\end{tikzpicture}
\end{center}
\end{frame}

% 3 directed
\begin{frame}{\hlorange{Đồ thị có hướng (Directed graph)}}
\begin{block}{\hlgreen{Định nghĩa}}
Trong đồ thị \hlred{có hướng}, mỗi cạnh có hướng: \(\hlblue{u} \rightarrow \hlblue{v}\).
\end{block}

\begin{itemize}
\item Dùng cho luồng, biểu đồ dependency.
\item Đỉnh có {\em \hlcyan{in-degree}} và {\em \hlcyan{out-degree}}.
\end{itemize}

\begin{center}
\begin{tikzpicture}[->,>=stealth']
\tikzset{v/.style={circle,draw=myred,fill=myred!18,minimum size=8mm}}
\node[v] (1) at (0,0) {\hlblue{1}};
\node[v] (2) at (2,1) {\hlblue{2}};
\node[v] (3) at (2,-1) {\hlblue{3}};
\draw (1) -- (2); \draw (2) -- (3); \draw (1) -- (3);
\end{tikzpicture}
\end{center}
\end{frame}

% 4 Representations - edge list, adj matrix, adj list
\begin{frame}[t]{\hlorange{Biểu diễn đồ thị} — \hlred{Danh sách cạnh} / \hlgreen{Ma trận kề} / \hlpurple{Danh sách kề}}
\scriptsize
\textbf{\hlorange{Danh sách cạnh (Edge list)}}: \\[2pt]
$\{(1,2,4),\ (1,4,2),\ (2,3,3),\ (4,3,1)\}$ \quad (\hlcyan{mỗi bộ: $(u,v,w)$})

\medskip
\textbf{\hlgreen{Ma trận kề (Adjacency matrix)} — 0/1:}
\begin{center}
\begin{tabular}{c|cccc}
\rowcolor{tablehead}\color{tableheadtxt} & \color{tableheadtxt}1 & \color{tableheadtxt}2 & \color{tableheadtxt}3 & \color{tableheadtxt}4\\\hline
1 & 0 & 1 & 0 & 1\\
2 & 1 & 0 & 1 & 0\\
3 & 0 & 1 & 0 & 1\\
4 & 1 & 0 & 1 & 0
\end{tabular}
\end{center}

\medskip
\textbf{\hlpurple{Danh sách kề (Adjacency list)} — hiệu quả khi đồ thị thưa:}
\begin{itemize}\scriptsize
\item \hlblue{1}: (2,4), (4,2)
\item \hlblue{2}: (1,4), (3,3)
\item \hlblue{3}: (2,3), (4,1)
\item \hlblue{4}: (1,2), (3,1)
\end{itemize}
\end{frame}

% 5 mapping from edge list to matrix & list (overlay highlight)
\begin{frame}{\hlorange{Danh sách cạnh} (minh hoạ)}
\small
Danh sách cạnh: \(\{(1,2,4),(1,4,2),(2,3,3),(4,3,1)\}\)

\medskip
\only<1>{
\begin{center}
\begin{tikzpicture}
\node at (0,0) {\begin{tabular}{c|cccc}
 & 1 & 2 & 3 & 4 \\\hline
1 & 0 &  &  &  \\
2 &  & 0 &  &  \\
3 &  &  & 0 &  \\
4 &  &  &  & 0
\end{tabular}};
\node at (4.5,0) {\begin{tabular}{l}
1:  \\
2:  \\
3:  \\
4:
\end{tabular}};
\end{tikzpicture}
\end{center}
}
\only<2>{
\begin{center}
\begin{tikzpicture}
\node at (0,0) {\begin{tabular}{c|cccc}
 & 1 & 2 & 3 & 4 \\\hline
1 & 0 & \cellcolor{mycyan!30}1 &  & \cellcolor{mycyan!30}1 \\
2 & \cellcolor{mycyan!30}1 & 0 &  &  \\
3 &  &  & 0 &  \\
4 & \cellcolor{mycyan!30}1 &  &  & 0
\end{tabular}};
\node at (4.5,0) {\begin{tabular}{l}
1: (\hlorange{2},4), (\hlorange{4},2)\\
2:  \\
3:  \\
4: 
\end{tabular}};
\end{tikzpicture}
\end{center}
}
\only<3>{
\begin{center}
\begin{tikzpicture}
\node at (0,0) {\begin{tabular}{c|cccc}
 & 1 & 2 & 3 & 4 \\\hline
1 & 0 & 1 &  & 1 \\
2 & 1 & 0 & \cellcolor{mycyan!30}1 &  \\
3 &  & \cellcolor{mycyan!30}1 & 0 & \cellcolor{mycyan!30}1 \\
4 & 1 &  & 1 & 0
\end{tabular}};
\node at (4.5,0) {\begin{tabular}{l}
1: (2,4), (4,2)\\
2: (1,4), (3,3)\\
3: (2,3), (4,1)\\
4: (1,2), (3,1)
\end{tabular}};
\end{tikzpicture}
\end{center}
}
\end{frame}

% 6 DFS theory
\begin{frame}{\hlorange{DFS} — Depth First Search}
\begin{block}{\hlgreen{Ý tưởng}}
Duyệt sâu nhất có thể trước khi quay lui (backtrack).
\end{block}
\medskip
\textbf{Thuộc tính}
\begin{itemize}
\item Dùng \hlpurple{stack} (hoặc \hlpurple{đệ quy}).
\item Độ phức tạp: \(\hlblue{O(V+E)}\).
\item Ứng dụng: kiểm tra liên thông, phát hiện chu trình, tạo cây DFS, phân tích thành phần liên thông.
\end{itemize}
\end{frame}

% 7 DFS pictorial step-by-step
\begin{frame}{\hlorange{Minh họa DFS} — thứ tự thăm (step-by-step)}
\centering
\begin{tikzpicture}[scale=1]
\tikzset{v/.style={circle,draw=myblue,fill=myblue!10,minimum size=8mm}}
\node[v] (1) at (0,2) {\hlblue{1}};
\node[v] (2) at (2,3) {\hlblue{2}};
\node[v] (3) at (4,2) {\hlblue{3}};
\node[v] (4) at (2,1) {\hlblue{4}};
\node[v] (5) at (4,0) {\hlblue{5}};
\draw (1)--(2) (2)--(3) (1)--(4) (2)--(4) (3)--(5) (4)--(5);
\only<1>{\node at (2,-1) {\hlcyan{Bắt đầu từ 1}};}
\only<2>{\node at (2,-1) {\hlcyan{Visited: 1}}; \node[fill=mycyan!40,circle,inner sep=1pt] at (0,2) {};}
\only<3>{\node at (2,-1) {\hlcyan{Visited: 1,2}}; \node[fill=mycyan!40,circle,inner sep=1pt] at (2,3) {};}
\only<4>{\node at (2,-1) {\hlcyan{Visited: 1,2,3}}; \node[fill=mycyan!40,circle,inner sep=1pt] at (4,2) {};}
\only<5>{\node at (2,-1) {\hlcyan{Visited: 1,2,3,5}}; \node[fill=mycyan!40,circle,inner sep=1pt] at (4,0) {};}
\only<6>{\node at (2,-1) {\hlcyan{Visited: 1,2,3,5,4 (kết thúc)}}; \node[fill=mycyan!40,circle,inner sep=1pt] at (2,1) {};}
\end{tikzpicture}
\end{frame}

% 8 DFS code
\begin{frame}[fragile]{\hlorange{Code DFS} (đệ quy)}
\begin{lstlisting}[language=C++]
vector<int> adj[N];
bool vis[N];

void dfs(int u){
  vis[u]=true;                  // danh dau
  // xu li u (in ra, do thoi gian
  for(int v: adj[u]){
    if(!vis[v]) dfs(v);         // di sau
  }
}
\end{lstlisting}
\end{frame}

% 9 BFS theory
\begin{frame}{\hlorange{BFS} — Breadth First Search}
\begin{block}{\hlgreen{Ý tưởng}}
Duyệt theo từng lớp (\hlpurple{level}) từ đỉnh bắt đầu; dùng \hlpurple{queue}.
\end{block}
\medskip
\textbf{Ứng dụng}: tìm đường ngắn nhất trong đồ thị không trọng số, khoảng cách theo số cạnh.
\end{frame}

% 10 BFS pictorial with levels
\begin{frame}{\hlorange{Minh họa BFS} — các tầng (levels)}
\begin{center}
\begin{tikzpicture}[scale=1]
\tikzset{v/.style={circle,draw=mygreen,fill=mygreen!12,minimum size=8mm}}
\node[v] (1) at (0,2) {\hlblue{1}};
\node[v] (2) at (-2,0) {\hlblue{2}};
\node[v] (3) at (0,0) {\hlblue{3}};
\node[v] (4) at (2,0) {\hlblue{4}};
\node[v] (5) at (0,-2) {\hlblue{5}};
\draw (1)--(2) (1)--(3) (1)--(4) (3)--(5);
\only<1>{\node at (0,-3) {\hlcyan{Bắt đầu từ 1 (level 0)}};}
\only<2>{ \node[fill=mycyan!40,rounded corners] at (0,2.8) {\hlblue{Level 0: 1}};}
\only<3>{
 \node[fill=mycyan!40,rounded corners] at (0,1.8) {\hlblue{Level 0: 1}};
 \node[fill=mycyan!20,rounded corners] at (-2,-1.2) {\hlgreen{Level 1: 2,3,4}};
}
\only<4>{
 \node[fill=mycyan!40,rounded corners] at (0,1.8) {\hlblue{Level 0: 1}};
 \node[fill=mycyan!20,rounded corners] at (0,-1.2) {\hlgreen{Level 2: 5}};
}
\end{tikzpicture}
\end{center}
\end{frame}

% 11 BFS code
\begin{frame}[fragile]{\hlorange{Code BFS}}
\begin{lstlisting}[language=C++]
vector<int> adj[N];
bool vis[N];

void bfs(int s){
  queue<int> q;
  q.push(s);
  vis[s]=true;
  while(!q.empty()){
    int u=q.front(); q.pop();
    // xử lý u
    for(int v:adj[u]){
      if(!vis[v]){
        vis[v]=true;
        q.push(v);
      }
    }
  }
}
\end{lstlisting}
\end{frame}

% 12 Dijkstra theory
\begin{frame}{\hlorange{Dijkstra} — Ý tưởng chính}
\begin{itemize}
\item Tìm đường ngắn nhất từ 1 đỉnh đến các đỉnh khác trong đồ thị có trọng số không âm.
\item Sử dụng \hlpurple{priority queue} (min-heap) để chọn đỉnh có \hlblue{dist} nhỏ nhất chưa xử lý.
\item Độ phức tạp: \(\hlblue{O((V+E)\log V)}\).
\end{itemize}
\end{frame}

% 13 Dijkstra example with table step-by-step
\begin{frame}{\hlorange{Dijkstra} — Ví dụ và bảng cập nhật}
\small
Đồ thị: cạnh (1,2,4),(1,4,2),(2,3,3),(4,3,1). Tìm từ s=\hlblue{1}.

\medskip
\only<1>{
\begin{tabular}{c|c}
\hlgreen{Đỉnh} & \hlgreen{dist (khởi tạo)}\\\hline
1 & 0\\
2 & $\infty$\\
3 & $\infty$\\
4 & $\infty$
\end{tabular}
}
\only<2>{
\begin{tabular}{c|c}
\hlgreen{Đỉnh} & \hlgreen{dist (sau relax cạnh từ 1)}\\\hline
1 & 0\\
2 & 4\\
3 & $\infty$\\
4 & 2
\end{tabular}
}
\only<3>{
\begin{tabular}{c|c}
\hlgreen{Đỉnh} & \hlgreen{dist (sau chọn 4 và relax)}\\\hline
1 & 0\\
2 & 4\\
3 & 3  % via 4: 2+1
\\
4 & 2
\end{tabular}
}
\only<4>{
\begin{tabular}{c|c}
\hlgreen{Đỉnh} & \hlgreen{dist (cuối)}\\\hline
1 & 0\\
2 & 4\\
3 & 3\\
4 & 2
\end{tabular}
}
\end{frame}

% 14 Dijkstra code
\begin{frame}[fragile]{\hlorange{Code Dijkstra} (mẫu)}
\begin{lstlisting}[language=C++]
using pii = pair<int,int>;
vector<pair<int,int>> adj[N];
vector<int> dist(N, 1e9);

void dijkstra(int s){
  priority_queue<pii, vector<pii>, greater<pii>> pq;
  dist[s]=0; pq.push({0,s});
  while(!pq.empty()){
    auto [d,u]=pq.top(); pq.pop();
    if(d>dist[u]) continue;
    for(auto [v,w]: adj[u]){
      if(dist[v] > dist[u] + w){
        dist[v] = dist[u] + w;
        pq.push({dist[v], v});
      }
    }
  }
}
\end{lstlisting}
\end{frame}

% 15 Flood Fill theory
\begin{frame}{\hlorange{Flood Fill} (Loang bảng)}
\begin{itemize}
\item Dùng DFS/BFS trên lưới để tô vùng liên thông.
\item Ứng dụng: tô màu trong paint, đếm vùng, giải mê cung.
\end{itemize}
\end{frame}

% 16 Flood Fill pictorial animation
\begin{frame}{\hlorange{Minh họa Flood Fill} — bước lan}
\centering
\begin{tikzpicture}[scale=0.7]
\foreach \x in {0,...,4}{
  \foreach \y in {0,...,4}{
    \draw (\x,\y) rectangle (\x+1,\y+1);
  }
}
\only<1>{\node at (2.5,-1) {\hlcyan{Bắt đầu tại ô (2,2)}};}
\only<2>{ \fill[mycyan!40] (2,2) rectangle (3,3); }
\only<3>{
 \fill[mycyan!40] (2,2) rectangle (3,3);
 \fill[mycyan!40] (1,2) rectangle (2,3);
 \fill[mycyan!40] (3,2) rectangle (4,3);
 \fill[mycyan!40] (2,3) rectangle (3,4);
 \fill[mycyan!40] (2,1) rectangle (3,2);
}
\only<4>{
 \node at (2.5,-1) {\hlgreen{Tất cả ô kề cùng màu được tô}};
 \fill[mycyan!40] (2,2) rectangle (3,3);
 \fill[mycyan!40] (1,2) rectangle (2,3);
 \fill[mycyan!40] (3,2) rectangle (4,3);
 \fill[mycyan!40] (2,3) rectangle (3,4);
 \fill[mycyan!40] (2,1) rectangle (3,2);
}
\end{tikzpicture}
\end{frame}

% 17 Kruskal theory
\begin{frame}{\hlorange{Kruskal} — Minimum Spanning Tree (MST)}
\begin{itemize}
\item Mục tiêu: chọn tập cạnh nối tất cả đỉnh với tổng trọng số nhỏ nhất.
\item Ý tưởng: sắp xếp cạnh theo trọng số tăng dần, chọn cạnh nếu không tạo chu trình (\hlpurple{DSU}).
\item Độ phức tạp: \(\hlblue{O(E\log E)}\).
\end{itemize}
\end{frame}

% 18 Kruskal pictorial step
\begin{frame}{\hlorange{Minh họa Kruskal} — chọn cạnh từng bước}
\centering
\begin{tikzpicture}[scale=0.9]
\tikzset{v/.style={circle,draw=myorange,fill=myorange!20,minimum size=8mm}}
\node[v] (1) at (0,0) {\hlblue{1}};
\node[v] (2) at (2,1.5) {\hlblue{2}};
\node[v] (3) at (4,0) {\hlblue{3}};
\node[v] (4) at (2,-1.5) {\hlblue{4}};
\only<1>{ \draw (4)--node[below]{\hlcyan{1}} (3); \draw (1)--node[left]{\hlcyan{2}} (4); \draw (2)--node[above]{\hlcyan{3}} (3); \draw (1)--node[above]{\hlcyan{4}} (2);}
\only<2>{ \draw[very thick,myred] (4)--node[below]{1} (3); \draw (1)--node[left]{2} (4); \draw (2)--node[above]{3} (3); \draw (1)--node[above]{4} (2); \node at (2,-2.5) {\hlgreen{Chọn (4,3) weight=1}};}
\only<3>{ \draw[very thick,myred] (4)--node[below]{1} (3); \draw[very thick,myred] (1)--node[left]{2} (4); \draw (2)--node[above]{3} (3); \draw (1)--node[above]{4} (2); \node at (2,-2.5) {\hlgreen{Chọn (1,4) weight=2}};}
\only<4>{ \draw[very thick,myred] (4)--node[below]{1} (3); \draw[very thick,myred] (1)--node[left]{2} (4); \draw[very thick,myred] (2)--node[above]{3} (3); \draw (1)--node[above]{4} (2); \node at (2,-2.5) {\hlgreen{Chọn (2,3) weight=3 — MST hoàn thành}};}
\end{tikzpicture}
\end{frame}

% 19 Kruskal code (DSU)
\begin{frame}[fragile]{\hlorange{Code Kruskal} (DSU)}
\begin{lstlisting}[language=C++]
struct Edge{int u,v,w;};
vector<Edge> E;
int parent[N], rnk[N];

int findp(int x){ return parent[x]==x?x:parent[x]=findp(parent[x]); }
void unite(int a,int b){
  a=findp(a); b=findp(b);
  if(a==b) return;
  if(rnk[a]<rnk[b]) swap(a,b);
  parent[b]=a;
  if(rnk[a]==rnk[b]) rnk[a]++;
}

sort(E.begin(),E.end(),[](Edge &a, Edge &b){return a.w<b.w;});
for(auto e:E){
  if(findp(e.u)!=findp(e.v)){
    unite(e.u,e.v);
    mst += e.w;
  }
}
\end{lstlisting}
\end{frame}

% 20 comparison table
\begin{frame}{\hlpurple{So sánh nhanh các thuật toán}}
\begin{center}
\begin{tabular}{l|c|c|c}
\rowcolor{tablehead}\color{tableheadtxt} Thuật toán & \color{tableheadtxt}Cấu trúc & \color{tableheadtxt}Dùng cho & \color{tableheadtxt}Độ phức tạp\\\hline
\hlblue{DFS} & stack/đệ quy & thành phần liên thông, chu trình & \(\hlblue{O(V+E)}\)\\
\hlgreen{BFS} & queue & đường ngắn không trọng số & \(\hlgreen{O(V+E)}\)\\
\hlorange{Dijkstra} & PQ & đường ngắn không âm & \(\hlorange{O(E\log V)}\)\\
\hlpurple{Kruskal} & DSU & MST & \(\hlpurple{O(E\log E)}\)
\end{tabular}
\end{center}
\end{frame}

% 21 exercises
\begin{frame}{\hlcyan{Bài tập \& hướng dẫn tự học}}
\begin{enumerate}
\item Viết \hlblue{DFS/BFS} in ra thứ tự duyệt của đồ thị mẫu.
\item Cài \hlorange{Dijkstra} cho đồ thị có trọng số trên và chạy từ đỉnh \hlblue{1}.
\item Cài \hlpurple{Kruskal} và kiểm tra MST.
\item Bài mở rộng: \hlred{Bellman–Ford}, \hlred{Floyd–Warshall}, \hlred{Topological sort}.
\end{enumerate}
\end{frame}

% 22 references / kết thúc
\begin{frame}{\hlgreen{Tài liệu tham khảo \& Kết luận}}
\begin{itemize}
\item Sách: \hlblue{Cormen et al.}, \hlblue{Kleinberg-Tardos}.
\item Trang: \hlpurple{CP-algorithms}, \hlpurple{GeeksforGeeks}, \hlpurple{TopCoder}.
\end{itemize}

\bigskip
\centerline{\hlorange{Cảm ơn — Hỏi \& Thảo luận}}
\end{frame}

\end{document}